
\phantomsection
\section*{Заключение}
\addcontentsline{toc}{section}{Заключение}

В результате выполнения курсового проекта была рассмотрена предметная область автоматизации частной медицинской клиники и определены основные проблемы, возникающие при ручной или разрозненной организации записи пациентов, расписания врачей и хранения медицинской информации.

В первой главе проведён анализ существующих программных решений. Сравнение было выполнено в формате «критерий - платформа»: для Инфоклиника.RU, Medesk, ONDOC, Cliniccards и разрабатываемого приложения рассмотрены онлайн-запись, личный кабинет пациента, электронная медицинская карта, управление расписанием, ролевая модель доступа, работа с документами, администрирование и открытость архитектуры. Анализ показал, что готовые системы обладают широкими возможностями, но являются закрытыми коммерческими продуктами, поэтому целесообразна разработка собственного веб-приложения.

Также в первой главе выполнен анализ существующих стеков разработки веб-приложений. Были рассмотрены варианты Python + Flask, Python + Django, Node.js + Express + React/Vue, PHP + Laravel, Java + Spring Boot и ASP.NET Core. По итогам сравнения выбран стек Python + Flask + PostgreSQL/SQLite + Bootstrap, так как он обеспечивает высокую скорость разработки, модульность и безопасность при создании CRUD-интерфейсов, форм записи и личных кабинетов.

Во второй главе выполнено проектирование веб-приложения. Были определены роли пользователей (пациент, врач, администратор), описаны пользовательские сценарии, сформулированы функциональные требования, разработаны ER-диаграмма, логическая и физическая схемы базы данных, а также подготовлены макеты страниц (Wireframes). Особое внимание уделено таблицам, связанным с пользователями, расписанием, записями на приём, медицинскими картами и документами.

В третьей главе приводится описание программной реализации разработанного веб-приложения на стеке Flask + SQLAlchemy + Bootstrap. В ходе работы были созданы все запланированные модули:
\begin{enumerate}
    \item Модуль авторизации и разграничения доступа на базе Flask-Login;
    \item Личный кабинет пациента с отображением электронной медицинской карты, списков будущих приёмов и истории посещений;
    \item Модуль онлайн-записи с динамической подгрузкой свободных слотов времени через асинхронные JavaScript-запросы к API;
    \item Рабочее место врача с расписанием на день, формой заполнения протоколов осмотра (жалобы, диагноз, рекомендации) и загрузкой файлов результатов исследований;
    \item Панель администратора для ведения справочников персонала, услуг и управления реестром записей на приём.
\end{enumerate}

Разработанная система контейнеризации с использованием Docker и Docker Compose обеспечивает простоту сборки и развертывания веб-приложения на любых серверах без необходимости предварительной ручной настройки СУБД и интерпретатора Python. Все функциональные требования были верифицированы в ходе тестирования, подтвердившего стабильность работы бизнес-логики и корректность отображения адаптивного интерфейса на экранах различных разрешений.

Таким образом, цель курсового проекта успешно достигнута: разработано и протестировано полнофункциональное веб-приложение для частной клиники, обеспечивающее надежное автоматизированное взаимодействие пациентов, медицинского персонала и администрации. Перспективами развития проекта являются интеграция системы SMS или Email-оповещений пациентов о записях, добавление платежного шлюза для онлайн-оплаты услуг, разработка API для интеграции с внешними лабораторными системами и усиление механизмов защиты персональных медицинских данных пациентов в соответствии с нормативными требованиями.
